% =====================================================================
%  SPLICE: The computability split -- first-use introduction
%  DATED 2026-08-09. Purpose: give the term a proper pedagogical
%  introduction at its first load-bearing appearance. Recommended
%  placement: opening of the Hopf-justification material, immediately
%  before the stabiliser computation of Proposition \ref{prop:gsm};
%  JS to confirm the seam. Sectioning level (\subsection) may need
%  adjusting to the host section.
%
%  External labels referenced (verify all present in the corpus):
%    sec:division-algebras, sec:self-reference, sec:fanout-boundary,
%    prop:gsm, prop:ghz, prop:w-access, lem:cproj,
%    app:setup, app:winding-number, lem:winding-number.
%
%  Citation keys used, already in the ledger:
%    baez2002octonions, furey2012unified, furey2018three,
%    krasnov2021spin9, duboisviolettetodorov2019exceptional,
%    szangolies2025standardmodel.
%  VERIFY-KEY (mint or confirm before compile; sourced from memory):
%    gunaydingursey1973  (Gunaydin & Gursey, J. Math. Phys. 14, 1651)
%    conwaysmith2003     (Conway & Smith, On Quaternions and
%                         Octonions, A K Peters, 2003, Sect. 8.4)
%
%  Assumes the amsthm environment `definition' exists in the preamble
%  alongside proposition/lemma; if not, add
%    \newtheorem{definition}[theorem]{Definition}
%  Pure ASCII throughout; British English; no bullets.
%  REV 2026-08-09b: carpet / self-referential-imposition imagery
%    installed after the sorting sentence, with sibling-instance
%    caution; strike freely if the paper voice should stay plainer.
% =====================================================================
% ===== BEGIN SPLICE: computability-split-first-introduction 2026-08-09 =====

\subsection{The computability split}
\label{sec:computability-split}

The octonions arrive from \S\ref{sec:division-algebras} with no marked
directions. The automorphism group $G_2 = \operatorname{Aut}(\mathbb{O})$
acts transitively on the unit imaginary octonions
\autocite{baez2002octonions}, so every imaginary direction is the equal of
every other; nothing in the algebra prefers one. The Standard Model gauge
group is nevertheless about to appear, in Proposition~\ref{prop:gsm}, as
the stabiliser of one of them. Between those two sentences sits the step on
which every algebraic route from $\mathbb{O}$ to $G_{\mathrm{SM}}$ turns,
and it deserves a more careful introduction than the literature usually
gives it. Three questions want answers before the term below can carry
weight: what exactly is being chosen, why some choice is unavoidable, and
what --- in the present framework --- does the choosing.

\begin{definition}[Computability split]
\label{def:computability-split}
Let $\ell$ be a unit imaginary octonion and $J = L_{\ell}$ left
multiplication by $\ell$, a complex structure on $\mathbb{O}$. The
\emph{computability split} is the $J$-invariant orthogonal decomposition
\begin{equation}
  \mathbb{O} \;\cong\; \mathbb{C}\ell \,\oplus\, \mathbb{C}^{3},
  \label{eq:computability-split}
\end{equation}
in which $\mathbb{C}\ell = \operatorname{span}_{\mathbb{R}}\{1,\ell\}$ is
called the \emph{record plane} and the six-dimensional complement,
organised by $J$ into the three planes cut out by the Fano lines through
$\ell$, is called the set of \emph{colour planes}. In the conventions of
Appendix~\ref{app:setup}, $\ell = e_1$, with colour planes
$\{e_2,e_3\}$, $\{e_4,e_5\}$ and $\{e_7,e_6\}$.
\end{definition}

The mathematics of the definition is classical: the stabiliser of $\ell$
in $G_2$ is $SU(3)$, acting on the colour planes as the fundamental
representation and on the record plane trivially, and that stabiliser
computation is how Proposition~\ref{prop:gsm} will extract
$G_{\mathrm{SM}}$. What is not classical is the name, and the name is the
point of this subsection.

The first two questions have the same answer across the division-algebra
literature: what is chosen is a preferred complex direction, and the
choice is unavoidable because the gauge group is its stabiliser --- no
selection, no subgroup. Furey selects an idempotent whose stabiliser is
this $SU(3)$ \autocite{furey2012unified,furey2018three}; Krasnov selects a
complex structure $J_R$ on the $Spin(9)$ spinor space $\mathbb{O}^2$
\autocite{krasnov2021spin9}; Dubois-Violette and Todorov select the
subalgebra $\mathbb{C} \subset \mathbb{O}$
\autocite{duboisviolettetodorov2019exceptional}; Szangolies selects the
direction that reduces the three-qubit geometry from $9{+}1$ to $3{+}1$
dimensions \autocite{szangolies2025standardmodel}; the tradition reaches
back to G\"unaydin and G\"ursey \autocite{gunaydingursey1973}. In every
case the selection is an input: the right structure emerges, but only by
assuming a preference the algebra itself does not contain.

The framework's answer to the third question is that no external hand is
needed, because the selection is made by any completed measurement. This
is a framework claim, not established mathematics, and we flag it as
such. In the language of \S\ref{sec:self-reference} and
\S\ref{sec:fanout-boundary}, an observer is a record-keeper: a system
that copies information into classical, readable, erasure-priced form. A
measurement, in the framework's geometric description, is the completion
of a closed information loop on $S^7$, and completing such a loop is
multiplication by a unit imaginary octonion. Which unit is immaterial in
a precise sense: choosing one determines the complementary six, since the
six left multiplications of the complement generate the Clifford algebra
$\mathbb{C}\mathrm{l}(6)$ whose volume element is $L_{\ell}$ itself
(Appendix~\ref{app:winding-number}) --- Conway and Smith's ``seven rights
make a left'' \autocite{conwaysmith2003}. Different measurements select
different $\ell$; all selections are related by $G_2$; and the
stabiliser, hence the gauge group, is the same for every one. The
derivation does not need a particular direction. It needs the fact of a
direction, and the fact of a direction is what ``a measurement has
occurred'' means. Within a fixed register the selection is pinned
further still: the orbit-wide cancellations behind the confinement
theorem hold only for the direction whose Fano blocks align with the
real--imaginary pairing of the register's amplitudes (the alignment
analysis following Lemma~\ref{lem:cproj}), so a register singles out its
own $\ell$ rather than borrowing an arbitrary one.

The name records what distinguishes the two summands relative to the
observer doing the selecting. A record is built sequentially ---
operations composed one after another --- and composition of operations
is associative whatever the operations are, so a record can transcribe
without loss only structure that is associatively organised. The record
plane is such a structure: a commutative associative subalgebra in which
a phase is a winding number and a winding number is, by
Lemma~\ref{lem:winding-number}, a ladder count --- the natural currency
of bookkeeping. It is the same plane in which the observer's own
amplitudes were assembled in \S\ref{sec:self-reference}, the analytic
continuation whose $2\pi\mathrm{i}$ ambiguities became windings. The
complement is nothing of the kind. It is not closed under multiplication
--- the colour planes multiply into one another and back into the record
plane along the Fano lines --- and triples drawn from distinct colour
planes feel the associator, so no single sequential transcript
coordinatises it. Its content is carried collectively or not at all,
which is the structural seed of confinement
(Proposition~\ref{prop:ghz}). The split is a \emph{computability} split
because it sorts the algebra by what a record can hold: the serialisable
part is the record plane, and the remainder is the address to which the
founding non-computability of \S\ref{sec:self-reference} is relocated.
The diagonal argument guarantees a residue that no self-description
captures; the split says where that residue lives.

% CAUTION (house rule: each instance carries its own proof): the
% paragraph below formalises the ALGEBRAIC instance of the
% bump-in-the-carpet principle. The dynamical instance --- relocation
% to the QM/GR interface at the FANOUT boundary, Paper 2's business ---
% is a sibling requiring its own statement, not a citation to this one.
This is the precise form of two heuristic images the framework has
carried from the start. The first is a carpet nailed flat in a room:
the room is what a record can hold, the carpet is the non-computability
the diagonal argument guarantees, and a carpet that cannot be removed
must bump somewhere. The split is the nailing, the colour planes are
the bump's address, and the transitivity of $G_2$ is the statement that
the bump can be relocated but never flattened --- every choice of
$\ell$ sorts the same algebra and strands the same excess. The second
image is the imposition of a non-self-referential structure on an
underlying self-referential one: a record is a transcript licensed to
forget bracketing --- associativity is exactly that licence --- while
the octonions' bracketing is content, and the associator is what the
imposition cannot flatten. The two images coincide because the
framework's founding inference runs from self-reference to
non-computability; the split gives them both their address.

The strata should be kept distinct. The
decomposition~\eqref{eq:computability-split}, the transitivity of $G_2$,
and the stabiliser computation are established mathematics. The
identification of the distinguished plane with the record plane, and of
measurement as the selecting event, are framework claims; their
consistency tests are the theorems the identification makes provable ---
the sector populations of Propositions~\ref{prop:ghz}
and~\ref{prop:w-access}, and the operator identity of
Lemma~\ref{lem:winding-number}, which would have no reason to hold if
the record plane and the algebra's selected plane were different planes.
% OPTIONAL (JS to decide; strike if the EWSB material stays in a later
% paper): The dynamical face of the selection --- its identification
% with electroweak symmetry breaking, the Higgs expectation value as the
% moment the direction freezes --- is developed in [TODO-REF].

% ===== END SPLICE: computability-split-first-introduction 2026-08-09 =====
